\uuid{GhSs}
\titre{Changements de variables}
\auteur{hanani ; mzari}
\organisation{exo7}
\licence{CC BY-NC-SA 4.0}
\datecreate{2020-01-20}
\niveau{L1}
\module{Analyse}
\chapitre{Calcul d'intégrales}
\sousChapitre{Autre}
\difficulte{3}
\theme{Calculs d'intégrales, changements de variables}
\exo7qcmid{lille-2.integrales.023}

\contenu{
\begin{question}
\qclassification{127.99}
\qtags{motcle=changements de variables}
Parmi les affirmations suivantes, cocher celles qui sont vraies :
\begin{answers}
\good{Le changement de variable $t=\sin x$ donne $\displaystyle \int _{\pi/6}^{\pi/4}\frac{\mathrm{d}\, x}{\sin x\cos x}=\int _{1/2}^{1/\sqrt{2}}\frac{\mathrm{d}\, t}{t(1-t^2)}$.}
\bad{$\forall t\in \Rr\setminus \{-1,0,1\}$, $\displaystyle \frac{1}{t(1-t^2)}=\frac{1}{t}+\frac{1}{1-t}-\frac{1}{1+t}$.}
\bad{Une primitive de $\displaystyle \frac{1}{t(1-t^2)}$ sur $]0,1[$ est $\displaystyle F(t)=\ln \frac{t}{1-t^2}$.}
\good{$\displaystyle \int _{\pi/6}^{\pi/4}\frac{\mathrm{d}\, x}{\sin x\cos x}=\ln \sqrt{3}$.}
\end{answers}
\vskip3mm
\begin{explanations}
Avec $t=\sin x\Rightarrow \mathrm{d}t=\cos x\, \mathrm{d}x$, on a : $t(\pi/6)=1/2$, $t(\pi/4)=1/\sqrt{2}$ et
$$\displaystyle \int _{\pi/6}^{\pi/4}\frac{\mathrm{d}\, x}{\sin x\cos x}=\int _{\pi/6}^{\pi/4}\frac{\cos x\, \mathrm{d}x}{\sin x\cos ^2x}=\int _{\pi/6}^{\pi/4}\frac{\cos x\, \mathrm{d}x}{\sin x(1-\sin ^2x)}=\int _{1/2}^{1/\sqrt{2}}\frac{\mathrm{d}\, t}{t(1-t^2)}.$$
Or, $\displaystyle \frac{1}{t(1-t^2)}=\frac{1}{t}+\frac{1}{2-2t}-\frac{1}{2+2t}$. Donc
$$\displaystyle \int _{1/2}^{1/\sqrt{2}}\frac{\mathrm{d}\, t}{t(1-t^2)}=\left[\ln \frac{t}{\sqrt{1-t^2}}\right]_{1/2}^{1/\sqrt{2}}=\ln \sqrt{3}.$$
\end{explanations}
\end{question}
}
